\section{Spectral curves and their classification}
\label{sec:spectral-geometry}

The Lax pencil associates a spectral curve to each trajectory. For
fixed matrix sizes $L,M,N$, the phase coefficients and initial modes
determine a family of marked curves. We describe its regular members,
their genus, and how special spectral levels accommodate curves of
lower genus.

This fixed-size question complements the BA construction of
\cref{sec:general-BA-realization}: any prescribed smooth curve admits
a compatible BA realization, enlarging the mode lists when needed.
The classification here describes how curves appear within particular
networks, while that construction chooses a network to accommodate
the curve.

\subsection{From the Lax pencil to a fixed curve}

For each spectral parameter $\nu$, consider the eigenvalue equation
\begin{equation}\label{eq:geometry-eigenvalue}
 \mathsf L(\nu,t)\psi=\lambda\psi,\qquad \psi\ne0.
\end{equation}
Its solvability condition is precisely
\begin{equation}\label{eq:geometry-affine}
 F(\lambda,\nu)=\det\bigl(\lambda\Id_d-\mathsf X(t)-\nu\mathsf D\bigr)=0,
 \qquad d=L+M+N.
\end{equation}
Thus the spectral curve records the eigenvalues as $\nu$ varies.
At a generic value of $\nu$ there are $d$ eigenvalues, so projection
to $\nu$ gives a degree-$d$ covering of the projective line.
Its branch values occur when eigenvalues coincide; they can be
located by eliminating $\lambda$ from $F=\partial_\lambda F=0$.

The equation is fixed along the Hamiltonian flow. Indeed,
\cref{eq:canonical-lax} implies $\partial_tF=0$, since
the trace of the adjugate of $\lambda\Id_d-\mathsf L$ multiplied
by $[\mathsf L,\mathsf T]$ vanishes. The phase coefficients and
initial modes determine the coefficients of $F$, which are the
spectral first integrals. Different integral values can therefore
give different curves, while a single trajectory preserves its curve.
This is the standard spectral-curve mechanism for matrix Lax systems
\cite{JRW1991,Beauville1990,AdlerVanMoerbeke2018}.

To obtain a compact smooth curve, first take the projective closure
of \cref{eq:geometry-affine}, then its normalization. Normalization
separates the local branches at singular points. We denote this smooth
curve by $C$; its genus, rather than the degree of the determinant,
is the genus relevant to the associated Abelian functions.

\subsection{Branches at infinity and genus}

Introduce a homogeneous coordinate $\tau$, with affine chart $\tau=1$.
The projective equation and its restriction to the line at infinity are
\begin{align}
 \widehat F(\lambda,\nu,\tau)
 &=\det\bigl(\lambda\Id_d-\nu\mathsf D-\tau\mathsf X\bigr)=0,
 \label{eq:geometry-projective}\\
 \widehat F(\lambda,\nu,0)
 &=(\lambda-\nu)^L(\lambda+\nu)^M\lambda^N.
 \label{eq:geometry-infinity}
\end{align}
The plane curve has degree $d$ and meets infinity at
$[1:1:0]$, $[-1:1:0]$ and $[0:1:0]$. The three points contain,
respectively, $L$, $M$ and $N$ branches. Their affine expansions are
\begin{equation}\label{eq:geometry-branches}
\begin{aligned}
 \lambda&=\nu-2b_\ell^L+O(\nu^{-1}),&&1\leq\ell\leq L,\\
 \lambda&=-\nu-2b_m^M+O(\nu^{-1}),&&1\leq m\leq M,\\
 \lambda&=b_n^N+O(\nu^{-1}),&&1\leq n\leq N.
\end{aligned}
\end{equation}
In particular, the phase coefficients specify the branch intercepts.
Distinct intercepts within each list give distinct tangent lines at
the corresponding plane point. Such a point is called an
\emph{ordinary multiple point}: all its branches are smooth and
their tangent lines are distinct. Multiplicity one means a smooth
point. Every branch in \cref{eq:geometry-branches} meets the infinity
line transversely.

We call a spectral curve \emph{regular} when the phase coefficients
are pairwise distinct within each of the three lists, its equation is
irreducible, and it has no affine singularities. Its only possible
plane singularities are then the prescribed ordinary points at infinity.

\begin{proposition}\label{prop:geometry-genus}
The normalization of a regular spectral curve has genus
\begin{equation}\label{eq:geometry-genus}
 G=LM+MN+NL-L-M-N+1.
\end{equation}
It has $d$ marked points at infinity, one for each branch in
\cref{eq:geometry-branches}.
\end{proposition}
\begin{proof}
For example, near $[1:1:0]$, use $\nu=1$ and write
$\lambda=1+\tau h$. The first nonzero coefficient in
$\widehat F(1+\tau h,1,\tau)$ is a nonzero constant times
\[
 \tau^L\prod_{\ell=1}^{L}(h+2b_\ell^L).
\]
Its distinct roots give the $L$ tangent lines. The other two points
give the factors
$\prod_{m=1}^{M}(h+2b_m^M)$ and
$\prod_{n=1}^{N}(h-b_n^N)$ in the same way.
An ordinary point of multiplicity $q$ reduces the genus of a degree-$d$
plane curve by $q(q-1)/2$. The plane genus formula
\cite[Chapter~IV, \S1]{Hartshorne1977} therefore gives
\[
 G=\frac{(d-1)(d-2)}2
   -\frac{L(L-1)+M(M-1)+N(N-1)}2,
\]
which simplifies to \cref{eq:geometry-genus}.
\end{proof}

This also explains the count in \cref{eq:canonical-genus-preview}:
the number of mode pairs is $p=G+(d-1)$, with $d-1$ independent
power charges. The spectral curve supplies the remaining geometric
part of that count.

The regular locus is nonempty when phases and initial modes vary;
Appendix~\ref{sec:geometry-details} proves this by smoothing a
union of lines.

\subsection{Regular curves at fixed matrix sizes}

The infinity branches provide a characterization in terms of marked
points and meromorphic functions. Write their divisors as
\begin{equation}\label{eq:geometry-markings}
 \mathcal D_L=\sum_{\ell=1}^{L}P_\ell^L,\qquad
 \mathcal D_M=\sum_{m=1}^{M}P_m^M,\qquad
 \mathcal D_N=\sum_{n=1}^{N}P_n^N.
\end{equation}
Here $P_\ell^L$, $P_m^M$ and $P_n^N$ are the points of $C$
corresponding to the infinity branches in \cref{eq:geometry-branches}.
Each point occurs once, and the three sets are disjoint.
For a meromorphic function $f$, the pole divisor $(f)_\infty$
lists its poles with their orders. Its degree is the number of
preimages of a generic value of $f$, counted with multiplicity.

\begin{theorem}\label{thm:geometry-classification}
Fix positive $L,M,N$, and let $G$ be given by
\cref{eq:geometry-genus}. For a smooth connected complex curve $C$
of genus $G$, the following statements are equivalent.
\begin{enumerate}
 \item $C$ is the normalization of a regular spectral curve of
 \cref{eq:canonical-pencil}, for some phase coefficients and initial modes.
 \item $C$ has a birational plane model of degree $d=L+M+N$
 with three distinct collinear ordinary points of multiplicities
 $L,M,N$, no other singularities, and with their common line
 transverse to every branch.
 \item $C$ has divisors as in \cref{eq:geometry-markings} and
 meromorphic functions $\lambda,\nu$ generating its function field
 $\C(C)$, with exact pole divisors
 \begin{equation}\label{eq:geometry-poles}
 \begin{aligned}
  (\lambda-\nu)_\infty&=\mathcal D_M+\mathcal D_N,\\
  (\lambda+\nu)_\infty&=\mathcal D_L+\mathcal D_N,\\
  (\lambda)_\infty&=\mathcal D_L+\mathcal D_M.
 \end{aligned}
 \end{equation}
\end{enumerate}
\end{theorem}
The equivalence is proved in Appendix~\ref{sec:geometry-details}.
The marked functions specify the regular curve family at these sizes.
Varying the phase coefficients and spectral integral values explores
this family; each Hamiltonian trajectory preserves its own spectral
curve. A prescribed curve of lower genus can occur at a special level,
as described below.

\subsection{Covering degrees and hyperelliptic cases}

The three cancellations in \cref{eq:geometry-poles} give maps to
the projective line of degrees
\begin{equation}\label{eq:geometry-degrees}
 \deg\lambda=L+M,\qquad
 \deg(\lambda-\nu)=M+N,\qquad
 \deg(\lambda+\nu)=L+N.
\end{equation}
Any two of these functions generate $\C(C)$, since they recover
$\lambda,\nu$ by linear combinations. These are concrete
restrictions on the curve, additional to its genus.

\begin{proposition}\label{prop:geometry-hyperelliptic}
A regular spectral curve of genus $G\geq2$ is hyperelliptic if and
only if at least two of $L,M,N$ equal one. The case $L=M=N=1$
has genus one.
\end{proposition}
Two singleton groups give a degree-two map in
\cref{eq:geometry-degrees}. Appendix~\ref{sec:geometry-details}
proves the converse using the covering degrees.
For example, in the regular $(1,M,1)$ case the degree-two function
$\lambda+\nu$ has pole divisor $Q_1+R_1$: its involution exchanges
these points, so the BA puncture is not a branch point. The
odd-degree specialization of \cref{sec:BA-realization} instead marks
a Weierstrass point as its single puncture. Its conditions $L,M\geq g$
come from interpolation at that marking, not from a restriction on
all hyperelliptic spectral curves. The general residue criterion
covers both settings; in the singleton group, the principal divisor
of $\lambda+\nu$ supplies the required relation between the two
holomorphic evaluation vectors.



\noindent\begin{minipage}{\textwidth}
For example, the following table lists the covering degrees from
\cref{eq:geometry-degrees}, in increasing order.
\begin{center}
\begin{tabular}{ccll}
 $(L,M,N)$ & $G$ & Covering degrees & Regular curve\\\hline
 $(1,1,1)$ & $1$ & $2,2,2$ & elliptic\\
 $(1,1,2)$ & $2$ & $2,3,3$ & hyperelliptic\\
 $(1,2,2)$ & $4$ & $3,3,4$ & non-hyperelliptic\\
 $(1,2,3)$ & $6$ & $3,4,5$ & non-hyperelliptic\\
 $(2,2,2)$ & $7$ & $4,4,4$ & non-hyperelliptic\\
 $(2,2,3)$ & $10$ & $4,5,5$ & non-hyperelliptic
\end{tabular}
\end{center}
\end{minipage}\par\medskip
In particular, the $(1,2,2)$ family is trigonal: it has a
degree-three map and no degree-two map. The $(2,2,2)$ family is
tetragonal. A hypothetical degree-three map and one of its
degree-four maps would generate the function field, since their
degrees are coprime; Castelnuovo--Severi would then give $G\leq6$.

There is also a useful product-surface model when $N=1$.
Using $\lambda-\nu$ and $\lambda+\nu$ as coordinates gives a
smooth curve of bidegree $(L+1,M+1)$ in
$\mathbb P^1\times\mathbb P^1$, with genus $LM$.
Here bidegree records the degrees in the first and second coordinates
of its defining equation; the corresponding projection degrees are
$M+1$ and $L+1$. For $N>1$, the same change gives bidegree
$(L+N,M+N)$, but its $N$ branches at the common infinity corner
share a tangent. Their total defect is $N(N-1)$: each pair has
intersection multiplicity two, as follows from their distinct
intercepts in \cref{eq:geometry-branches}. Consequently its genus is
\[
 (L+N-1)(M+N-1)-N(N-1)=G.
\]

\subsection{A symmetric reduction and its Prym}

A symmetric reduction with $L=M=s$, $N=1$ gives an involution
$\iota(\lambda,\nu)=(\lambda,-\nu)$. The moving divisor then
lies in a translate of its Prym, the connected part of the Jacobian
reversed by $\iota$. On a regular level the curve has genus $s^2$,
and the Prym has dimension $s(s+1)/2$.
Appendix~\ref{sec:geometry-prym-proof} gives the mode constraints,
the dimension count and the eigenvector-divisor proof.

\subsection{Special spectral levels}

Keep the phase coefficients distinct within each group and suppose
the plane curve remains irreducible. If particular integral values
create affine singularities, its normalization has genus
\begin{equation}\label{eq:geometry-special-genus}
 g(C)=G-\sum_{P\in\operatorname{Sing}_{\mathrm{aff}}(F)}\delta_P.
\end{equation}
The sum runs over all affine singular points; $\delta_P$ is their
local genus defect, equal to one at an ordinary node.
Thus the sizes determine the regular genus, while the spectral
integrals determine whether a particular level has lower genus.
Coincident phase coefficients can also change the infinity
singularities and require their defects to be recomputed.

For birational prescribed BA realizations,
\cref{prop:BA-prescribed-spectral-curve} identifies the prescribed smooth
curve with the normalization and gives the total excess defect $G-g$.
For example, the $(L,M,N)=(3,3,1)$ network has regular genus nine,
whereas the genus-three trigonal realization in
\cref{sec:numerical-validation} has excess defect six. The additional
singularities belong to its plane spectral image; the prescribed BA
curve is its smooth normalization.
